% !PRINTED 134
\weberpage{134}
%<<CONT
genau $(p-1)\nu$ Elemente vorkommen, deren Grad durch $p$ theilbar ist, ist also nur noch die Formel

\[ \varepsilon)\quad \nu = j\,r_2 \]

zu beweisen. Diese ergiebt sich aus folgender Ueberlegung. Nach $\alpha')$ ist die Anzahl der Elemente $U$ gleich $p\,\nu$. Unter diesen Elementen $U$ ist aber gewiss eines, das Einheitselement, dessen Grad nicht durch $p$ theilbar ist, und folglich ist nach $\delta)$

\[ j\,r_2\,(p-1) < p\,\nu. \]

Andererseits ist $n = \mu\,\nu = p\,r_1 r_2 j$, und da $\nu$ relativ prim zu $p\,r_1$ ist, so ist $j\,r_2$ durch $\nu$ theilbar. Also ist $j\,r_2 : \nu$ eine ganze positive Zahl, die kleiner als $p : p-1$ und folglich auch kleiner als 2 ist, und die daher nur gleich 1 sein kann. Dadurch ist $\varepsilon)$ bewiesen und zugleich das ganze Theorem.

\begin{center}
\S.\,33.

\gesperrt{Gruppen vom Grade $p^{\alpha} q$.}
\end{center}

\gesperrt{Frobenius} hat in der erwähnten Abhandlung den Satz ausgesprochen:

VI. \gesperrt{Jede Gruppe vom Grade $p^{\alpha} q$ ist, wenn $p$ und $q$ von einander verschiedene Primzahlen sind und $\alpha$ einen beliebigen positiven Exponenten bedeutet, metacyklisch.}

Der folgende Beweis dieses Satzes entstammt einer brieflichen Mittheilung von \gesperrt{Frobenius} an den Verfasser.

Zunächst ist klar, dass es genügt, zu beweisen, dass keine Gruppe $P$ vom Grade $p^{\alpha} q$ einfach ist. Unser Satz ist nämlich bewiesen für $\alpha = 1$. Nehmen wir ihn also für jedes kleinere $\alpha$ schon als bewiesen an, und setzen voraus, dass $P$ einen echten Normaltheiler $Q$ habe, der mehr als das Einheitselement umfasst, so sind die beiden Gruppen $Q$, $P/Q$ nach der Voraussetzung oder nach dem Satze IV metacyklisch, und also ist auch $P$ metacyklisch.

Wir nehmen also jetzt an, es sei die Gruppe $P$ vom Grade $p^{\alpha} q$ einfach, und wir haben aus dieser Annahme einen Widerspruch abzuleiten, um ihre Unmöglichkeit nachzuweisen.

$\alpha)$ Zunächst haben wir nach dem Satze I., \S.\,29 einen Theiler von $P$ vom Grade $q$. Wir wollen annehmen, es gebe einen
%CONT>>
